\DocumentMetadata{
  lang        = en-US,
  tagging     = on,
}

\RequirePackage{shellesc}
\documentclass[11pt, oneside]{article}   	% use "amsart" instead of "article" for AMSLaTeX format

\usepackage{geometry}                		% See geometry.pdf to learn the layout options. There are lots.
\geometry{letterpaper}                   		% ... or a4paper or a5paper or ... 
%\geometry{landscape}                		% Activate for rotated page geometry
\usepackage[parfill]{parskip}    		% Activate to begin paragraphs with an empty line rather than an indent
\usepackage{graphicx}				% Use pdf, png, jpg, or eps§ with pdflatex; use eps in DVI mode
								% TeX will automatically convert eps --> pdf in pdflatex		
\usepackage{amssymb}
\usepackage{amsmath}
\usepackage{tablefootnote}
\usepackage{multirow}
\usepackage{tabularx}
\renewcommand\tabularxcolumn[1]{m{#1}}

\usepackage{dsfont}
\usepackage[hang,flushmargin]{footmisc}
\usepackage[normalem]{ulem}

\usepackage{pgfplots,pgfplotstable}

\usepackage[autosize,forceshell]{dot2texi}
\usepackage{tikz}
\usetikzlibrary{shapes,arrows,patterns}
\usepackage{tkz-fct}

\ExplSyntaxOn
\cs_new_protected:Npn \suspendtagging #1 { \tag_suspend:n{#1} }
\cs_new_protected:Npn \resumetagging  #1 { \tag_resume:n{#1} }
\ExplSyntaxOff

\makeatletter
\long\def\ifnodedefined#1#2#3{%
    \@ifundefined{pgf@sh@ns@#1}{#3}{#2}%
}

\pgfplotsset{
    discontinuous/.style={
    scatter,
    scatter/@pre marker code/.code={
        \ifnodedefined{marker}{
            \pgfpointdiff{\pgfpointanchor{marker}{center}}%
             {\pgfpoint{0}{0}}%
             \ifdim\pgf@y>0pt
                \tikzset{options/.style={mark=*}}
                \draw [dashed,thin] (marker-|0,0) -- (0,0);
                \draw plot [mark=*,mark options={fill=white}] coordinates {(marker-|0,0)};
             \else
                \tikzset{options/.style={mark=none}}
             \fi
        }{
            \tikzset{options/.style={mark=none}}        
        }
        \coordinate (marker) at (0,0);
        \begin{scope}[options]
    },
    scatter/@post marker code/.code={\end{scope}}
    }
}

\makeatother

\newcommand*\Eval[3]{\left(#1\middle)\right\rvert_{#2}^{#3}}


\usepackage{arydshln}
\usepackage{mathtools}


%\usepackage{tagpdf}
%\tagpdfsetup{activate,paratagging,interwordspace}
\usepackage{accsupp}
\usepackage{hyperref}
\usepackage{lipsum}
%\usepackage{axessibility}
%\usepackage{accessibility}

%SetFonts
\DeclareEmphSequence{\bfseries\itshape}
\DeclareMathOperator {\ints}{\mathbb{Z}}
\DeclareMathOperator {\reals}{\mathbb{R}}
\DeclareMathOperator*{\argmin}{argmin}
\DeclareMathOperator*{\argmax}{argmax}
\DeclareMathOperator{\E}{\mathbb{E}}
\DeclareMathOperator{\V}{\mathbb{V}}


\renewcommand\labelenumi{(\theenumi)}
\renewcommand*{\thefootnote}{\fnsymbol{footnote}}

\usepackage{tcolorbox}

\newenvironment{aside}[1][Aside]{\begin{tcolorbox}[colback=black!5,colframe=black!75!black,title=#1]}{\end{tcolorbox}}


\title{Review via Linear Regression}
\author{02-680: Essentials of Mathematics and Statistics}
%\date{}							% Activate to display a given date or no date

\begin{document}
\maketitle

\section{Linear Regression}
The \emph{linear regression} problem is as follows 
(we covered the least squares solution briefly in Topic 12 (Calculus), written there as $Ax=b$ with $x$ as the unknown; here the unknown is $\beta$ and $x$ is the data):
\begin{itemize}
\item[\textbf{$x$}] --- Independent variable[s] (what we control)
\item[\textbf{$y$}] --- Dependent variable[s] (the measured values)
\item[\textbf{$\beta_s$}] -- slope
\item[\textbf{$\beta_0$}] -- intercept
\end{itemize} 


Some examples of regression in biology:
\begin{center}\begin{tabular}{|c|c|}
\hline
Independent & Dependent\\
\hline\hline
Expression of genes expressed in immune response & asthma severity\\
\hline
Expression of regulator genes & Expression of target genes\\
\hline
Clinical variables & Insulin level\\
\hline
... & ...
\end{tabular}\end{center}

Let's assume we have some $x,y\in\reals^n$, for $n$ samples where each $(x_i,y_i)$ are corresponding. 
Then for each sample $y_i = \beta_s x_i + \beta_0$, or for all samples at once \[y = \beta_s x + \beta_0\mathbf{1},\]
where $\mathbf{1}\in\reals^n$ is the vector of all ones.

To simplify things let's define:
$X = \begin{bmatrix} x_1 & 1\\x_2 & 1\\\vdots & \vdots\\x_n & 1\end{bmatrix}$ and $\beta = \begin{bmatrix}\beta_s\\\beta_0\end{bmatrix}$, 
then the equation above becomes simply:
\[y = X\beta.\]

Recall (possibly) from earlier that we want to find $\beta$ that minimizes the $\|y-X\beta\|^2_2$, 
or the square of the $L_2$-norm. 

\begin{align*}
\|y-X\beta\|^2_2 & = (y-X\beta)^T(y-X\beta)\\
& =  (y^T-\beta^TX^T)(y-X\beta)\\
& = y^Ty + \beta^TX^TX\beta - 2 \beta^TX^Ty
\end{align*}

To find the critical point of the curve (i.e. the minimum) we need
\begin{align*}
\frac{d}{d\beta} y^Ty + \beta^TX^TX\beta - 2 \beta^TX^Ty &= 0\\
\frac{d}{d\beta}\beta^TX^TX\beta  - \frac{d}{d\beta}2 \beta^TX^Ty &= 0\\
2 X^TX\beta - 2 X^Ty &=0\\
2 X^TX\beta &= 2 X^Ty \\
\beta &= (X^TX)^{-1}X^Ty
\end{align*}

\begin{aside}[When does this solution exist?]
The last step needs $X^TX$ to be invertible, and eigenelements (Topic 11) tell us when it is.
\begin{itemize}
\item $X^TX$ is symmetric, so its eigenvalues are real.
\item A matrix is singular exactly when its determinant is $0$, and the determinant is the product of the eigenvalues. So $X^TX$ is invertible if and only if \emph{none} of its eigenvalues is $0$.
\item A zero eigenvalue happens when the columns of $X$ are linearly dependent (Topic 10). For example, if every $x_i$ is the same value then the first column of $X$ is a multiple of the column of $1$s. With several predictors, the same thing happens if one is a combination of the others (e.g., one measurement is the sum of two others). Then $X^TX\beta=X^Ty$ has many solutions and $(X^TX)^{-1}$ does not exist.
\item For any eigenvector $v$ of $X^TX$ with eigenvalue $\lambda$, we have $\lambda\|v\|_2^2=v^TX^TXv=\|Xv\|_2^2\ge0$, so every eigenvalue is non-negative. This is why the critical point we found is a minimum of the squared error and not a maximum.
\end{itemize}
\end{aside}

What if I want to know the probability of some $y_i$ given $x_i$? (so $p(y_i\mid x_i)$)

Assuming it's some natural process, we can assume 
\[p(y_i\mid x_i) \sim \mathcal{N}(\beta_0 + \beta_s x_i, \sigma^2).\]
or equivalently
\[y_i = \beta_0 + \beta_s x_i + \epsilon_i, \epsilon_i \sim \mathcal{N}(0,\sigma^2).\]

So given the data, we can actually use MLE to find $\beta$!
Let $X_i=\begin{bmatrix}x_i & 1\end{bmatrix}$ be the $i$th row of $X$, so the mean of $y_i$ is $X_i\beta=\beta_0+\beta_s x_i$.

\begin{align*}
\ln p(\mathcal{D}\mid \beta, \sigma) &= \ln \left[\prod_{i=1}^{n} \frac{1}{\sigma\sqrt{2\pi}} e^{-\frac{1}{2\sigma^2}(y_i-X_i\beta)^2}\right]\\
	&= \sum_{i=1}^{n} \ln \left[\frac{1}{\sigma\sqrt{2\pi}} e^{-\frac{1}{2\sigma^2}(y_i-X_i\beta)^2}\right]\\
	&= \frac{-n}{2}\ln \sigma^2 - \frac{n}{2}\ln(2\pi) + \sum_{i=1}^{n} \ln e^{-\frac{1}{2\sigma^2}(y_i-X_i\beta)^2}\\
	&= \frac{-n}{2}\ln \sigma^2 - \frac{n}{2}\ln(2\pi) -\frac{1}{2\sigma^2}\sum_{i=1}^{n} (y_i-X_i\beta)^2
\end{align*}

\begin{align*}
\frac{\partial}{\partial\beta} {\color{red}\frac{-n}{2}\ln \sigma^2 - \frac{n}{2}\ln(2\pi)} -\frac{1}{2\sigma^2}\sum_{i=1}^{n} (y_i-X_i\beta)^2 &= \frac{\partial}{\partial\beta} -\frac{1}{2\sigma^2}\sum_{i=1}^{n} (y_i-X_i\beta)^2\\
&=  \frac{\partial}{\partial\beta} -\frac{1}{2\sigma^2} \|y-X\beta\|_2^2\\
&= -\frac{1}{2\sigma^2}  \frac{\partial}{\partial\beta} (y^Ty-2y^TX\beta+\beta^TX^TX\beta)\\
&= ... \\
\end{align*}

$\hat\beta_{MLE} = (X^TX)^{-1}X^Ty$, that's just Least Squares!

But we don't know $\sigma$, that's actually, what we want to know because it models how uncertain we are with any given output.

\begin{align*}
\frac{\partial}{\partial\sigma^2} \frac{-n}{2}\ln \sigma^2 {\color{red}- \frac{n}{2}\ln(2\pi)} -\frac{1}{2\sigma^2}\sum_{i=1}^{n} (y_i-X_i\beta)^2 &= \frac{-n}{2\sigma^2} +\frac{1}{2}(\sigma^2)^{-2}\sum_{i=1}^{n} (y_i-X_i\beta)^2
\end{align*}

$\hat\sigma^2_{MLE} = \frac{\sum_{i=1}^{n} (y_i-X_i\hat\beta_{MLE})^2}{n}$

With these we can determine for some new $\tilde y$ and $\tilde x$ what $p(\tilde y \mid \tilde x)$ is. 

\end{document}